% ============================================================
% DRAFT NOTE (problem)
% Contains: A plain-language statement of the forecasting target, then the
% formal estimand in three steps: the switch position, the paired
% continuations (fork and re-prefill variants), and the quality function.
% Written now: The complete settled estimand and label definition.
% Pending: None, apart from notation changes required by a venue template.
% Sources: docs/methodology.md; docs/goal.md; docs/_why/6_intervention_semantics.md;
% docs/_why/3_measuring_quality.md
% ============================================================
\section{Problem: Damage at a Switch Position}

Consider a model partway through generating an answer. The serving system
could activate KV-cache compression now and save memory for the rest of the
generation, but the finished answer might lose task quality as a result. The
quantity we study is exactly that cost: how much final quality the completed
answer loses if compression is activated at this moment, compared with
leaving the cache untouched. Defining this quantity requires three
ingredients: where the intervention happens, which two continuations are
compared, and how their quality is scored.

\paragraph{Where: the switch position.}
A switch position \switchpos{} means that exactly \switchpos{} generated
tokens are fixed and the intervention occurs immediately before the next
token is computed. We write \decoderstate{} for the complete decoder state at
this boundary: the KV cache together with everything else required to
continue generation exactly, such as pending input, position indices, and the
decoding, stopping, and numerical configuration. The full state matters
because at this boundary the most recent token may not yet be incorporated
into the cache.

\paragraph{What is compared: two paired continuations.}
From \decoderstate{} we create two independent, identical forks and compress
exactly one of them with compressor \compressor{} at removal ratio \ratio{}.
Both forks then generate to completion under the same prompt, emitted prefix,
token budget, EOS rules, decoding configuration, attention backend, and
numerical precision. Call the two completed outputs
\begin{equation}
y^{\mathrm{keep}} = \operatorname{continue}(\decoderstate),
\qquad
y^{\mathrm{comp}} = \operatorname{continue}\!\left(\mathcal{C}_{c,r}(\decoderstate)\right).
\label{eq:fork-arms}
\end{equation}
Damage is the signed difference in their final quality,
\begin{equation}
\dfork = q\!\left(y^{\mathrm{keep}}\right) - q\!\left(y^{\mathrm{comp}}\right).
\label{eq:fork-damage}
\end{equation}
Because compression is the only difference between the branches, the score
difference is attributable to activating compression at \switchpos.

Some compressors are not defined as a transformation of a live cache; their
faithful execution semantics is a compressed re-prefill of the prompt and
prefix. For those compressors both arms are rebuilt through the same no-press
re-prefill path $R$, so the control arm matches the treatment arm's execution
route, and the target is
\begin{equation}
\drefill = q\!\left(y_R^{\mathrm{keep}}\right) - q\!\left(y_R^{\mathrm{comp}}\right),
\label{eq:refill-damage}
\end{equation}
where $y_R^{\mathrm{keep}}$ and $y_R^{\mathrm{comp}}$ are built exactly as in
Equation~\ref{eq:fork-arms}, but from the re-prefilled state
$R(\decoderstate)$ instead of \decoderstate{} itself.
Section~\ref{sec:intervention} specifies which compressors receive which
target and how both are validated.

\paragraph{How it is scored: task-grounded quality.}
The initial task is IFEval. For a completed output $y$, quality is its
instruction-level loose accuracy,
\begin{equation}
q(y) = \frac{\text{instructions satisfied by } y}
{\text{instructions in the prompt}}.
\label{eq:quality}
\end{equation}
Damage is task grounded: different tokens or wording do not count as damage
unless the final score is lower. The label \damage{} is signed, with positive
values meaning degradation and negative values meaning compression-associated
lift, and we keep negative values during training. Positive harm
$\max(0,\damage)$ and any-damage and major-damage risks are derived reporting
quantities, not substitutes for the signed magnitude target. Strict IFEval
accuracy is retained as a secondary robustness score.
